\paragraph{Linear equations for the two completion columns.}
Work over \(k=\overline{\mathbb F}_5\). In the kernel chart write \(B\)
for the nonzero kernel direction with first star entry one and \(C,E\)
for the two completion columns with star entries \((0,1,0)\) and
\((0,0,1)\). Let \(q_j\) be the homogeneous term of degree \(j\) in the
permanent expanded about \(J_4\). The three pure equations on the five
free kernel edges are \(q_2(B)=q_3(B)=q_4(B)=0\). For each completion
column the terms linear in that column give the same polar map
\[
  T_B:X\longmapsto
    \bigl((dq_2)_B X,(dq_3)_B X,(dq_4)_B X\bigr).
\]
Its matrix has nine tangent columns. Exact elimination over
\(\mathbb F_5\) gives the unit ideal from the pure equations and the
maximal minors, proving rank three on the normalized pure kernel locus.
In a full level plane \(T_B\) kills the plane. Every tangent vector
can have its star entries canceled by a unique element of that plane,
since the star projection there is invertible. Hence the restriction
of \(T_B\) to the six coordinates with zero star entries still has
rank three at every possible completion. On any invertible minor,
three entries of \(C\) and three of \(E\) are determined by linear
algebra. This leaves eleven coordinates before localization, five
for \(B\) and three for each completion column. The three pure and
22 remaining coefficient equations still have to hold. Of the twenty
free-column minors, nineteen are nonzero polynomials. Modulo the pure
equations the ideal generated by all nineteen is already generated by
the three minors on column triples
\[
 (12,22,23),\qquad (11,12,22),\qquad (12,13,33).
\]
The two ideals with the pure equations adjoined have the same reduced
Gr\"obner basis. Thus these three open charts cover every completion.
Adjoining one inverse to each chart gives twelve variables and
26 equations on that chart.

\begin{proposition}[Exclusion of rational kernel lines at five]
\label{prop:rational-kernel-exclusion}
No geometric level three-plane in the collision branch at five has a
corrected-diagonal kernel line defined over \(\mathbb F_5\).
Consequently no complex level three-plane has a corrected-diagonal
kernel line defined over \(\mathbb Q\).
\end{proposition}
\begin{proof}
In edge order \(01,02,03,12,13,23\), test the projective kernel
directions for \(q_2=q_3=q_4=0\). Exact enumeration gives 24 points
of \(\mathbb P^5(\mathbb F_5)\) forming one \(S_4\)-orbit, represented
by \(b=(1,0,2,3,1,2)\). The test uses the permanent along each
direction at five scalars, which determines a polynomial of degree
at most four. It suffices to exclude one companion of this
representative with star \((0,1,0)\).

Write \(u=(C_{11},C_{12},C_{13})\), \(v=(C_{22},C_{23},C_{33})\).
The three polar rows have the form \(A u+L v+c=0\), with row degrees
one, two and three in the kernel edges. Put \(d=\det A\). In the
companion replace \(u\) by \(-\operatorname{adj}(A)(Lv+c)\), \(v\)
by \(dv\), and the fixed star by \(d(0,1,0)\). Denote the resulting
column by \(\widetilde C(B,v)\). At every actual companion it equals
\(dC\), including when \(d=0\). Each entry is homogeneous of degree
six in the kernel edges.

In \(\operatorname{per}(J_4+sB+t\widetilde C)\) take the six
coefficients of \(s^it^j\) with \(j\ge2,\ i+j\le4\), and their
monomial multiples in \(v\) through degree four, allowing multipliers
of degree at most \(4-j\). They give a \(39\times35\) matrix in the
monomial basis of \(v\). For a fixed direction this matrix \(M\) gives
an affine test \(My=0,\ y_{\boldsymbol 0}=1\). At an actual companion
the vector of evaluated monomials solves it. Thus inconsistency of
this test suffices to exclude a companion.
Every maximal minor vanishes at an extendible direction. The executable fixes 35
rows with determinant \(H\) of kernel degree 507, obtained by summing
their degrees \(i+6j\). At \(b\) it checks \(d=3,\ H(b)=2\) modulo
five and reconstructs a degree-four polynomial identity for one.
This certifies exclusion even for companions over \(k\).
Moreover \(H\) and its vertex-permutation translates vanish on all
geometric kernel directions of full planes. These equations apply
also where the chosen polar minor vanishes. In particular the
homogeneous \(S_4\)-invariant
\[
 H_{\mathrm{inv}}(b')=\sum_{\pi\in S_4}H(\pi b')
\]
vanishes there. The checker evaluates the summands on the 24 raw
relabelings of \(b\), before any projective rescaling, and obtains
\(H_{\mathrm{inv}}(b)=1\) in \(\mathbb F_5\). Its nonvanishing at
the representative excludes the full orbit.

For the complex conclusion, scale a rational kernel vector to have
primitive edge entries at five. Its diagonal entries are determined
by the corrected-kernel equations with denominators invertible at five.
It thus gives a line over \(\mathbb Z_{(5)}\) with a well-defined
\(\mathbb F_5\)-rational specialization. The locus of level planes
containing this section is closed in the projective Grassmannian.
A complex point would put the generic point in its proper image,
forcing a geometric point over five through the specified line.
The exclusion just proved rules this out.
\end{proof}
The remaining closed-branch search concerns kernel lines over
extensions of \(\mathbb F_5\) satisfying these determinantal equations
and all the completion equations. The special-fiber enumeration in
the proof concerns directions of the kernel. The plane containing such
a direction was allowed to live over the algebraic closure throughout.
